% AIPL / ABCL^c+ 言語ガイド
%   lualatex -interaction=nonstopmode aipl_guide_ja.tex   （2回）
\documentclass[11pt,a4paper]{ltjsarticle}
\usepackage{amsmath,amssymb}
\usepackage{listings}
\usepackage{xcolor}
\usepackage{geometry}
\usepackage{array}
\usepackage{booktabs}
\usepackage[hidelinks]{hyperref}
\geometry{margin=22mm}

\definecolor{codebg}{gray}{0.96}
\lstset{
  basicstyle=\ttfamily\small,
  backgroundcolor=\color{codebg},
  frame=single,
  framesep=4pt,
  rulecolor=\color{gray},
  breaklines=true,
  showstringspaces=false,
  columns=fullflexible,
  keepspaces=true,
  keywordstyle=\bfseries,
  commentstyle=\itshape\color{gray!130},
}
\lstdefinelanguage{AIPL}{
  morekeywords={class,method,var,float,new,now,future,await,reply,send,become,
    call,if,else,while,do,select,case,timeout,self,sender,remote,print,
    int,string,bool,unit,any},
  sensitive=true,
  morecomment=[l]{//},
  morecomment=[s]{/*}{*/},
  morestring=[b]",
}
\lstset{language=AIPL}

\newcommand{\ty}[1]{\ensuremath{\mathtt{#1}}}

\title{AIPL / ABCL$^{c}$+ 言語ガイド\\
{\large ---\ 型システムの改修を反映した現状の仕様と使い方 ---}}
\author{ABCL$^{c}$+ / AIOS 処理系（\texttt{yaskodama/abclcp}, branch \texttt{make-src-base}）}
\date{2026年07月30日 07:22 JST 版\\{\small 対象コミット: \texttt{04892cf}（branch \texttt{make-src-base}）}}

\begin{document}
\maketitle

\begin{abstract}
本書は AIPL（ABCL$^{c}$+）の\textbf{現状の}言語ガイドである。
第 \ref{part:change} 部で今回の改修による変更点と移行方法をまとめ、
第 \ref{part:lang} 部で言語仕様、第 \ref{part:sample} 部でサンプルプログラムを
ソースコードと実行結果つきで説明し、第 \ref{part:tool} 部でビルドと実行の手順を示す。

今回の改修の中心は、メソッドの戻り値型を \texttt{reply} から推論できるようにし、
さらに省略可能な戻り値型注釈 \texttt{method m(x) : T} を導入したことである。
これにより \texttt{now} / \texttt{future} の結果が \ty{any} ではなく
具体的な型を持つようになった。あわせて文字列連結を \texttt{+} から
\texttt{++} へ分離し、曖昧なオーバーロードをエラーに格上げし、
メソッド本体の仮引数をシグネチャに結線した。
設計の根拠と測定結果は \texttt{docs/aipl\_reply\_inference\_ja.pdf} にある。
\end{abstract}

\tableofcontents

\part{今回の変更点}
\label{part:change}

\section{変更の要約}

\begin{center}
\small
\begin{tabular}{@{}lp{7.6cm}l@{}}
\toprule
 & 内容 & 影響 \\
\midrule
戻り値型の推論 & \texttt{reply(v)} から各メソッドの戻り値型を推論する。
  \texttt{now} / \texttt{future} の結果が \ty{any} でなくなった & 追加のみ \\
戻り値型注釈 & \texttt{method m(x) : T} を書けるようになった（省略可） & 追加のみ \\
被覆検査 & \ty{unit} 以外を宣言したメソッドは全実行パスで \texttt{reply}
  しなければならない & 注釈時のみ \\
\texttt{reply} の線形性 & \texttt{reply} は高々一度。二重 \texttt{reply} は型エラー & \textbf{要確認} \\
select の照合 & \texttt{case m(...)} を \texttt{m} の署名に照合する（arity と型） & \textbf{要確認} \\
select の \texttt{reply} & case 本体の \texttt{reply} は\emph{選択されたメッセージ}への返信
  として型付けされる & 意味の明確化 \\
公開時の注釈必須 & \texttt{web\_expose} したアクターのメソッドは注釈が必須 & \textbf{要移行} \\
文字列連結 & \texttt{+} から \texttt{++} へ分離。\texttt{+} は数値専用 & \textbf{要移行} \\
整数演算 & \texttt{7 + 2} が \texttt{9.} ではなく \texttt{9} を返す。
  \texttt{/} は \ty{float} & 値が変わる \\
曖昧な overload & エラーに格上げ（従来は警告して既定候補を選択） & \textbf{要移行} \\
仮引数の結線 & 本体の仮引数がシグネチャの引数型と繋がった。
  引数の型検査が実際に効くようになった & \textbf{要確認} \\
\bottomrule
\end{tabular}
\end{center}

\section{移行が必要な点}

\subsection{文字列連結は \texttt{++} に書き換える}

\texttt{+} は数値専用になった。文字列を混ぜると型エラーになる。

\begin{lstlisting}
print("count = " + n);      // 変更前
print("count = " ++ n);     // 変更後
\end{lstlisting}

\texttt{++} は両辺を文字列化して連結するので、どちらが文字列でなくてもよい。
結合は \texttt{+} より\emph{弱い}ので、\texttt{"x=" ++ a + b} は
\texttt{"x=" ++ (a + b)} と読まれる。

機械的に移行するには、型検査器のエラー位置に従うのが安全である。
\texttt{no overload of + matches} が出た \texttt{(line, col)} は
演算子トークンそのものを指すので、その 1 文字を \texttt{++} にして
再検査すればよい。ただし列は\textbf{バイトオフセット}なので、
日本語を含む行ではバイト単位で置換する必要がある。

\subsection{公開するアクターのメソッドに注釈を付ける}

\texttt{web\_expose} で名指ししたアクターのメソッド（\texttt{init} を除く）は
戻り値型の注釈が必須になった。相手からは本体が見えないので、
返信の型は宣言でしか伝わらないためである。

\begin{lstlisting}
method say(msg) : string { reply("echo: " ++ msg); }
\end{lstlisting}

\texttt{web\_listen} だけを呼んでいて \texttt{web\_expose} していないアクターは
警告にとどまる（\texttt{POST /api/send} でアクター名を指定すれば到達できるが、
公開の意図があるとは限らないため）。
\texttt{AIOS\_LAX\_EXPOSE=1} でエラーを警告に落とせる。

\subsection{曖昧な式に注釈を付ける}

両辺が未束縛の \texttt{a + b} は、\ty{int} と \ty{float} のどちらとも
解釈できて principal type を持たない。従来は警告して \ty{float} を選んでいたが、
それは静的な型と実行結果が食い違う原因になっていたのでエラーにした。

\begin{lstlisting}
method add(a, b) { reply(a + b); }          // エラー: ambiguous overload
method add(a, b) : int { reply(a + b); }     // 注釈で確定する
\end{lstlisting}

宣言した型は\emph{期待型}として式の推論へ流れるので、注釈を書くだけで
オーバーロードが一意に決まる。\texttt{AIOS\_LAX\_OVERLOAD=1} で
従来の警告動作に戻せる。

\subsection{\ty{int} と \ty{float} の混在を直す}

仮引数の結線により、本体が要求する型と呼び出し側が渡す型が照合される
ようになった。そのため次のような不整合が表に出る。

\begin{lstlisting}
class Hello {
  float count = 0.;
  method init(n) { count = n; }   // 本体は n : float を要求する
}
var h = new Hello(5);             // int を渡している -> 型エラー
var h = new Hello(5.);            // 正しい
\end{lstlisting}

\subsection{二重 \texttt{reply} を直す}

\texttt{reply} は高々一度でなければならない。2 度目の \texttt{reply} は
待ち手がすでに 1 度目の値を受け取ったあとに来るので黙って捨てられていた。

\begin{lstlisting}
method m(k) : int { reply(k); reply(k + 1); }   // 型エラー
method m(k) : int { if (k > 0) { reply(1); } else { reply(2); } }  // OK
\end{lstlisting}

ループ本体の \texttt{reply} も「2 回以上ありうる」と判定される。

\part{言語ガイド}
\label{part:lang}

\section{プログラムの構成}

プログラムはクラス宣言とグローバル文の並びである。

\begin{lstlisting}
class Counter {
  var n = 0;                    // フィールド
  method init(start) {          // new のときに自動で呼ばれる
    n = start;
  }
  method bump() : int {
    n = n + 1;
    reply(n);
  }
}

var c = new Counter(10);        // グローバル文
print(now c.bump());
\end{lstlisting}

フィールドは \texttt{var x = e;} または \texttt{float x = e;} で宣言する。
メソッドは \texttt{method m(p1, ..., pn) \{ ... \}}、
戻り値型を書く場合は \texttt{method m(p1, ..., pn) : T \{ ... \}} である。

\texttt{init} という名前のメソッドは \texttt{new C(args)} のときに
自動で呼ばれる。引数の個数が合わないと生成が飛ばされる。

\section{型}

\begin{center}
\begin{tabular}{@{}ll@{}}
\toprule
型 & 説明 \\
\midrule
\ty{int}, \ty{float}, \ty{string}, \ty{bool}, \ty{unit} & 基底型 \\
\ty{T[]}          & 配列 \\
\ty{future\ T}    & \texttt{future} 送信の結果。\texttt{await} で \ty{T} を取り出す \\
\ty{actor(C)}     & クラス \texttt{C} のアクター。\texttt{new C(...)} の型 \\
\ty{any}          & 静的に内容を保証しない境界（\texttt{sender}、リモート送信先） \\
\bottomrule
\end{tabular}
\end{center}

戻り値型注釈に書けるのは
\texttt{int} / \texttt{float} / \texttt{string} / \texttt{bool} /
\texttt{unit} / \texttt{any} と、大文字で始まるクラス名（アクター型）である。

\begin{quote}
\small
メソッドの引数は\textbf{単相}である。シグネチャの引数型変数は全呼び出し地点で
共有されるので、同じメソッドを \ty{int} と \ty{string} で呼ぶと衝突する。
\end{quote}

\section{式}

\begin{center}
\begin{tabular}{@{}ll@{}}
\toprule
式 & 意味 \\
\midrule
\texttt{123}, \texttt{1.5}, \texttt{"abc"} & リテラル \\
\texttt{x} & 変数・フィールド・仮引数 \\
\texttt{a ++ b} & 文字列連結（両辺を文字列化）。結合は最も弱い \\
\texttt{a + b}, \texttt{a - b}, \texttt{a * b} & 数値演算。\ty{int} 同士は \ty{int}、
  \ty{float} 同士は \ty{float} \\
\texttt{a / b} & 除算。\ty{int} 同士でも \ty{float}（整数除算はしない） \\
\texttt{a > b}, \texttt{a < b}, \texttt{a >= b}, \texttt{a <= b} & 比較。結果は \ty{bool} \\
\texttt{new C(args)} & アクターの生成 \\
\texttt{now t.m(args)} & 同期送信。返信が来るまで待ち、その値になる \\
\texttt{future t.m(args)} & 非同期送信。\ty{future\ T} を返す \\
\texttt{await e} & future の解決を待つ \\
\texttt{f(args)} & 組込み関数の呼び出し \\
\texttt{( e )} & 括弧 \\
\bottomrule
\end{tabular}
\end{center}

送信先 \texttt{t} は変数名か \texttt{remote("host:port", "actorName")} である。

\begin{quote}
\small
\textbf{既知の欠落}：\texttt{==} と \texttt{!=} は字句解析器と評価器には
あるが、\emph{式の文法規則が無い}ため書けない。等値比較が必要な場合は
現状 \texttt{>=} と \texttt{<=} を組み合わせるしかない。
\end{quote}

\section{文}

\begin{center}
\begin{tabular}{@{}ll@{}}
\toprule
文 & 意味 \\
\midrule
\texttt{var x = e;} & 変数宣言 \\
\texttt{var x = new C(args);} & アクター生成つき宣言 \\
\texttt{x = e;} & 代入 \\
\texttt{f(args);} \quad \texttt{call f(args);} & 組込み関数の呼び出し \\
\texttt{send t.m(args);} & 非同期送信。返信を待たない \\
\texttt{send self.m(args);} & 自分自身への非同期送信 \\
\texttt{send sender.m(args);} & 送信元への非同期送信 \\
\texttt{send! t.m(args);} & 検査を緩めた送信 \\
\texttt{if (e) s} \quad \texttt{if (e) s else s} & 条件分岐 \\
\texttt{while e do s} & 繰り返し（括弧は不要） \\
\texttt{\{ s ... \}} & ブロック \\
\texttt{become C(args);} & 振る舞いの置換 \\
\texttt{select \{ ... \}} & メッセージの選択受信 \\
\bottomrule
\end{tabular}
\end{center}

\begin{quote}
\small
\texttt{now self.m()} と \texttt{now sender.m()} は\emph{文法上書けない}。
自分への \texttt{now} は必ず自己デッドロックになるので、書けないことが
そのまま安全側に働いている。
\end{quote}

\section{アクターと通信}

\subsection{三つの送信}

\begin{center}
\begin{tabular}{@{}lll@{}}
\toprule
 & 待つか & 式か文か \\
\midrule
\texttt{send a.m(x);}      & 待たない & 文 \\
\texttt{now a.m(x)}        & 返信まで待つ & 式（結果は宣言／推論された型） \\
\texttt{future a.m(x)}     & 待たない & 式（\ty{future\ T}。\texttt{await} で取る） \\
\bottomrule
\end{tabular}
\end{center}

\subsection{\texttt{reply}}

メソッドは値を「返す」のではなく \texttt{reply(v)} で呼び出し側の future を
解決する。型システムはメソッドごとに戻り値型を持ち、本体のすべての
\texttt{reply} と、すべての \texttt{now} / \texttt{future} 送信地点が
この型を共有する。

\begin{itemize}
\item \texttt{reply} は\textbf{高々一度}。
\item 戻り値型を \ty{unit} 以外で\emph{宣言}した場合は、
  \textbf{全実行パスでちょうど一度}になる。
\item \texttt{reply} がまったく無いメソッドは戻り値 \ty{unit} と扱われる。
  \texttt{send} で呼ぶ「通知」型のメソッドがこれにあたる。
\end{itemize}

\subsection{\texttt{select}}

\begin{lstlisting}
select {
  case job(k) -> { reply("done:" ++ k); }
  timeout 500 -> { print("timed out"); }
}
\end{lstlisting}

\texttt{case m(x1, ..., xn)} は同名メソッドの署名に照合される。
引数の個数と型が合わないと型エラーになる。

\textbf{case 本体の \texttt{reply} は、囲むメソッドではなく
「選択されたメッセージ」への返信である。}
処理系が case 実行の前に返信先を選択されたメッセージへ切り替えるためで、
上の例の \texttt{reply} は \texttt{job} への返信になる。したがって
\texttt{select} を含むメソッド自身は \ty{unit} でよい。
一方 \texttt{timeout} 本体は囲むメソッドの返信先のまま実行される。

\begin{quote}
\small
\texttt{select} は\emph{まだ処理されていない}メッセージを mailbox から探す。
通常のメソッド dispatch が先に同名メッセージを取ってしまうと
case には残らないので、\texttt{select} で待つ名前は直接呼ばせない設計に
するのが安全である（第 \ref{part:sample} 部 g4 の実行結果を参照）。
\end{quote}

\section{型検査が見るもの}

\begin{center}
\small
\begin{tabular}{@{}lp{9.0cm}@{}}
\toprule
検査 & 内容 \\
\midrule
メソッドの存在   & ローカル送信先に、その名前のメソッドがあるか \\
引数の個数と型   & 呼び出し側の実引数と、\emph{本体が要求する型}が一致するか \\
\texttt{reply} の型 & すべての \texttt{reply} が同じ（宣言された）型か \\
\texttt{reply} の回数 & 高々一度か。注釈があれば全パスでちょうど一度か \\
select の照合    & \texttt{case} の arity と引数型が署名と一致するか \\
公開時の注釈     & \texttt{web\_expose} したアクターのメソッドに注釈があるか \\
オーバーロード   & 候補が一意に決まるか（曖昧ならエラー） \\
\bottomrule
\end{tabular}
\end{center}

代表的なエラーメッセージを挙げる。

\begin{lstlisting}[language={}]
no method foo in actor(C)
type mismatch                                  <- 引数や代入の型が合わない
no overload of + matches (string, int)         <- + は数値専用
ambiguous overload of + for ('a, 'a): candidate return types are float / int
reply type mismatch: this method is declared to reply with int,
  but replies with string here
method m is declared to reply with int, but some execution path does not reply
method m may reply more than once on some path
select: case m binds 1 variable(s) but method m takes 2
actor C is reachable from outside this program, but m has no return type
  annotation; ...
\end{lstlisting}

\section{外部公開}

\texttt{web\_listen(port)} で HTTP ゲートウェイが起動し、
\texttt{web\_expose(path, actorName)} で公開エンドポイントに別名が付く。

\begin{center}
\begin{tabular}{@{}ll@{}}
\toprule
エンドポイント & 内容 \\
\midrule
\texttt{GET /}                & 簡易 UI \\
\texttt{POST /api/send}       & \texttt{to=<actor>\&method=<m>\&args=<a,b>} で任意のアクターへ \\
\texttt{POST /api/x/<key>}    & \texttt{web\_expose} した別名へ \\
\bottomrule
\end{tabular}
\end{center}

\texttt{POST /api/send} は\textbf{任意のアクターに名前で到達できる}。
すなわち境界を開いているのは \texttt{web\_listen} であって、
\texttt{web\_expose} は別名を付けるだけである。型検査はこれを踏まえ、
名指しされたアクターはエラー、それ以外は警告としている。

HTTP から渡される引数は文字列を \ty{int} $\to$ \ty{float} $\to$ \ty{string}
の順に解釈して作られる。境界が動的であること自体は注釈では消せないが、
\emph{どの型を約束しているか}は宣言できる。

\section{組込み関数（抜粋）}

\begin{center}
\small
\begin{tabular}{@{}ll@{}}
\toprule
分類 & 名前 \\
\midrule
基本         & \texttt{print}, \texttt{typeof}, \texttt{wait}, \texttt{reply}, \texttt{spawn} \\
数学         & \texttt{sin}, \texttt{cos}, \texttt{tan}, \texttt{asin}, \texttt{acos}, \texttt{atan},
               \texttt{sqrt}, \texttt{exp}, \texttt{log10}, \texttt{abs},
               \texttt{floor}, \texttt{ceil}, \texttt{round} \\
配列         & \texttt{array\_empty}, \texttt{array\_push}, \texttt{array\_get},
               \texttt{array\_set}, \texttt{array\_len} \\
Web          & \texttt{web\_listen}, \texttt{web\_expose} \\
描画 (SDL)   & \texttt{sdl\_init}, \texttt{sdl\_line}, \texttt{sdl\_erase\_line},
               \texttt{sdl\_clear}, \texttt{sdl\_present} \\
AI 呼び出し  & \texttt{ai\_call}, \texttt{ai\_call\_with\_system},
               \texttt{model\_generate}, \texttt{gemini\_generate}, \texttt{openai\_generate} \\
AIOS 内省    & \texttt{capabilities}, \texttt{aios\_kernel}, \texttt{aios\_actors},
               \texttt{aios\_actor\_info}, \texttt{aios\_mailbox\_len} \\
AIOS 通信    & \texttt{aios\_now}, \texttt{aios\_future}, \texttt{aios\_emit},
               \texttt{aios\_events} \\
記憶・タスク & \texttt{aios\_memory\_put}, \texttt{aios\_memory\_get},
               \texttt{aios\_task\_create}, \texttt{aios\_task\_set}, \texttt{aios\_tasks} \\
プロトコル   & \texttt{protocol\_define}, \texttt{protocol\_start},
               \texttt{protocol\_use}, \texttt{protocol\_state}, \texttt{protocol\_end} \\
リモート査読 & \texttt{remote\_review}, \texttt{remote\_review\_ja},
               \texttt{remote\_reviewer\_host} \\
\bottomrule
\end{tabular}
\end{center}

\part{サンプルプログラム}
\label{part:sample}

本節のサンプルは \texttt{docs/samples/guide/} にある。
すべて型検査を通り、実行結果も掲載のとおりである。

\section{g1: クラス・フィールド・\texttt{send}・文字列連結}

\lstinputlisting{samples/guide/g1_hello.abcl}

\begin{lstlisting}[language={}]
hello, AIPL
tick 1
tick 2
\end{lstlisting}

\texttt{init} は \texttt{new Greeter("AIPL")} のときに自動で呼ばれる。
\texttt{tick} は \texttt{reply} しないので \ty{unit} を宣言してあり、
被覆検査は課されない。\texttt{send} は返信を待たないので、
2 つの \texttt{tick} は順に処理される。
\texttt{count + 1} は \ty{int} 同士なので \ty{int} のままで、
表示も \texttt{1}, \texttt{2} と整数になる。

\section{g2: \texttt{now} / \texttt{future} / \texttt{await} と注釈}

\lstinputlisting{samples/guide/g2_now_future.abcl}

\begin{lstlisting}[language={}]
twice = 43
awaited = 20
v=7
\end{lstlisting}

\texttt{twice} は \ty{int} を宣言しているので \texttt{now c.twice(21)} は
\ty{int} になり、\texttt{s + 1} という算術に使える。
これは改修前にはできなかった（\texttt{now} の結果が \ty{any} だったため
\texttt{no overload of + matches (any, int)} で落ちていた）。

\texttt{label} は注釈を省いているが、\texttt{reply("v=" ++ a)} から
戻り値 \ty{string} が推論される。

\section{g3: 複数アクターと、アクターを跨いだ \texttt{now}}

\lstinputlisting{samples/guide/g3_actors.abcl}

\begin{lstlisting}[language={}]
ok, left=7
sold out
\end{lstlisting}

\texttt{Stock.take} が \ty{int} を返し、それが \texttt{Front.place} の中の
\texttt{now} の型になり、さらに \texttt{place} 自身の \ty{string} へ繋がる。
型はアクター境界を越えて伝播する。

\texttt{place} は \ty{string} を宣言しているので、
\texttt{if} の両方の枝で \texttt{reply} していることが検査される。
片方を消すと
\texttt{some execution path does not reply} になる。

なお \texttt{n = n - k} で \texttt{n} が \ty{int} なので \texttt{k} も
\ty{int} に決まり、呼び出し側の \texttt{place(3)} と整合する。
ここで \texttt{place("x")} と書けば\emph{引数の}型エラーになる
（改修前は素通りしていた）。

\section{g4: \texttt{select} と \texttt{reply} の帰属}

\lstinputlisting{samples/guide/g4_select.abcl}

\begin{lstlisting}[language={}]
direct:1
waiting
timed out
\end{lstlisting}

\texttt{case job(k)} の本体にある \texttt{reply} は \texttt{serve} ではなく
\texttt{job} への返信なので、\texttt{serve} は \ty{unit} のままでよい。
型検査器も case 本体では \texttt{reply} を \texttt{job} の戻り値型に
照合する。

実行結果は\textbf{この例が抱える設計上の注意}を示している。
\texttt{now w.job(1)} が通常のメソッド dispatch で先に処理されたため、
\texttt{serve} の \texttt{select} が動いたときには mailbox に
\texttt{job} が残っておらず、\texttt{timeout} に落ちた。
\texttt{select} で待つメッセージ名は、直接呼ばせない設計にするのが安全である。

\section{g5: 外部公開と注釈の必須化}

\lstinputlisting{samples/guide/g5_web.abcl}

このプログラムは型検査を通る。\texttt{say} の注釈を外すと次のようになる。

\begin{lstlisting}[language={}]
[Type error] line 16, col 1: actor Echo is reachable from outside this
  program, but say has no return type annotation; a remote caller cannot
  see the body, so the reply type has to be declared (method say(...) : T)
\end{lstlisting}

\texttt{web\_listen(8080)} を実行するとゲートウェイが起動して待ち受けに
入るので、確認は次のようにする。

\begin{lstlisting}[language={}]
$ curl -s -X POST http://localhost:8080/api/x/echo \
       -d 'method=say&args=hi'
\end{lstlisting}

\section{負例の一覧}

\texttt{docs/samples/reply\_inference/} には、\emph{通らないこと}を
確認するためのサンプルが揃っている。

\begin{center}
\small
\begin{tabular}{@{}ll@{}}
\toprule
ファイル & 何が起きるか \\
\midrule
\texttt{s2\_missing\_reply}      & \texttt{reply} 忘れの結果を算術に使う \\
\texttt{s3\_conflict}            & 分岐ごとに異なる型を \texttt{reply} \\
\texttt{s5\_overload\_ambiguity} & principal type が無い \texttt{a + b} \\
\texttt{s8\_annotation\_conflict}& 宣言と食い違う \texttt{reply} \\
\texttt{s9\_missing\_path}       & 一部のパスで \texttt{reply} しない \\
\texttt{s10\_expose\_unannotated}& 公開アクターに注釈が無い \\
\texttt{s12\_service\_exposed}   & 同上（統合サンプルの公開版） \\
\texttt{s14\_select\_pattern}    & \texttt{case} の arity が署名と違う \\
\texttt{s15\_double\_reply}      & 二重 \texttt{reply} \\
\bottomrule
\end{tabular}
\end{center}

\part{ビルドと実行}
\label{part:tool}

\section{処理系のビルド}

\begin{lstlisting}[language={}]
cd src && make thread_repl
\end{lstlisting}

\texttt{dune} でのビルドは \texttt{src/lexer.ml} の二重生成で失敗するため、
\texttt{src/Makefile} を使う。

\section{プログラムの実行}

REPL コマンドを書いたファイルを \texttt{-f} に渡す。

\begin{lstlisting}[language={}]
cat > run.repl <<'EOF'
load docs/samples/guide/g2_now_future.abcl
compile
EOF
./src/abclrepl_thread -q -f run.repl < /dev/null
\end{lstlisting}

\texttt{-f} に \texttt{.abcl} を直接渡すと REPL コマンドとして 1 行ずつ
解釈されてクラス定義が読まれないので、\texttt{load} と \texttt{compile} を
使うこと。REPL の主なコマンドは
\texttt{load} / \texttt{compile} / \texttt{list} / \texttt{send} /
\texttt{ast} / \texttt{pprint} / \texttt{script} / \texttt{exit} である。

\section{型検査だけを走らせる}

REPL は SDL に依存するので、型検査専用の小さなドライバを用意している。

\begin{lstlisting}[language={}]
ocamlfind ocamlc -package unix -thread -linkpkg -w -a -I src -o tc \
  src/location.cmo src/types.cmo src/ast.cmo src/typing_env.cmo \
  src/infer.cmo src/typecheck.cmo src/parser.cmo src/lexer.cmo \
  docs/samples/reply_inference/tc_main.ml

./tc docs/samples/guide/g3_actors.abcl
\end{lstlisting}

推論・宣言された戻り値型の表とクラスごとのシグネチャが表示される。

\section{型と実行値の差分テスト}

型検査が付けた戻り値型と、実行時に実際に \texttt{reply} された値の型を
突き合わせるハーネスがある。

\begin{lstlisting}[language={}]
python3 scripts/type_runtime_diff.py abclc/*.abcl docs/samples/guide/*.abcl
\end{lstlisting}

処理系側は \texttt{AIOS\_TYPE\_TRACE=1} のとき \texttt{reply} のたびに
\texttt{[rtype] Class\#method = tag} を出す。この仕組みで、整数演算が
\ty{float} を返していた不整合が実際に見つかった。
わざと食い違う負例には
\texttt{// TYPE\_RUNTIME\_DIFF: expect-mismatch} と書いておけば既知として扱われる。

\section{環境変数}

\begin{center}
\begin{tabular}{@{}ll@{}}
\toprule
変数 & 効果 \\
\midrule
\texttt{AIOS\_QUIET=1}          & 型スキーマ等のデバッグ出力を抑制 \\
\texttt{AIOS\_LAX\_OVERLOAD=1}  & 曖昧なオーバーロードをエラーではなく警告にする \\
\texttt{AIOS\_LAX\_EXPOSE=1}    & 公開アクターの注釈漏れをエラーではなく警告にする \\
\texttt{AIOS\_TYPE\_TRACE=1}    & \texttt{reply} のたびに \texttt{[rtype]} 行を出す \\
\texttt{AIOS\_MODEL\_PROVIDER}  & AI 呼び出しのプロバイダ選択 \\
\texttt{ABCL\_AI\_PROVIDER}     & 同上（旧名） \\
\texttt{REMOTE\_REVIEWER\_HOSTPORT} & \texttt{remote\_review} の宛先 \\
\bottomrule
\end{tabular}
\end{center}

\section{関連文書}

\begin{itemize}
\item \texttt{docs/aipl\_reply\_inference\_ja.pdf} ---
  今回の改修の設計根拠と測定結果。推論だけでは決まらない 6 点、
  既存アクター言語との比較、実装の詳細、退行試験の数値。
\item \texttt{docs/ABCLCP\_TYPE\_SOUNDNESS\_REPORT.md} --- 型健全性の整理。
\item \texttt{docs/samples/reply\_inference/README.md} --- 負例サンプルの一覧と意図。
\item \texttt{docs/AIOS\_LANGUAGE\_DESIGN.md}, \texttt{docs/CAPABILITIES.md} ---
  AIOS としての設計。
\end{itemize}

\section{現状の限界}

\begin{itemize}
\item \texttt{==} と \texttt{!=} が式として書けない（文法規則が無い）。
\item \texttt{sender} は \ty{any} で、\texttt{send sender.m(...)} は
  メソッドの存在も型も検査されない。
\item リモート送信先（\texttt{remote(...)}）の戻り値は \ty{any} のまま。
  対向ノードのインタフェース記述を読み込む仕組みが無い。
\item アクター型どちらも単一化で区別されない
  （\texttt{actor(A)} と \texttt{actor(B)} が単一化に成功する）。
\item \texttt{become C(...)} はコンストラクタの引数しか検査せず、
  置換後のインタフェースが元と一致するかを見ていない。
\item メソッドの引数は単相で、多相なメソッドは書けない。
\item \ty{int} と \ty{float} の間に暗黙変換が無く、
  両辺が未束縛の \texttt{a + b} は注釈が必要。
\end{itemize}

\end{document}
